\section{Proyeksi dan Subruang Terkomplemen}
\begin{defn}
 Pemetaan linear terbatas dari suatu ruang Banach ke dirinya sendiri disebut
 \index{proyeksi!di ruang Banach}%
 \index{operator!proyeksi}%
\df{proyeksi} jika bersifat idempoten. Jika $E$ adalah pemetaan seperti itu, kita menyebutnya
proyeksi \df{pada $\ran E$ sepanjang $\ker E$}.
\end{defn}

Perhatikan bahwa ini sudah pasti \emph{bukan} hal yang sama dengan proyeksi di
ruang Hilbert. Proyeksi ruang Hilbert adalah \emph{proyeksi ortogonal}; yakni,
proyeksi itu swaadjoin sekaligus idempoten. Jangkauan suatu proyeksi ruang Hilbert
tegak lurus terhadap kernelnya. Karena istilah ``ortogonal'', ``swaadjoin'', dan
``tegak lurus'' tidak bermakna di ruang Banach, kita tidak akan mengharapkan kedua pengertian
``proyeksi'' itu identik. Namun, dalam kedua kasus tersebut proyeksi \emph{memang}
terbatas, linear, dan idempoten.

\begin{defn} Misalkan $M$ subruang linear tertutup dari suatu ruang Banach~$B$. Jika terdapat
subruang linear tertutup $N$ dari $B$ sedemikian sehingga $B = M \oplus N$, maka kita mengatakan bahwa $M$ adalah
 \index{subruang terkomplemen}%
 \index{subruang!terkomplemen}%
subruang \df{terkomplemen}, bahwa $N$ adalah
 \index{komplemen!dari subruang Banach}%
\df{komplemen (ruang Banach)}-nya, dan bahwa subruang $M$ dan $N$ bersifat
\df{komplementer}.
\end{defn}

\begin{prop}\label{C069814} Misalkan $M$ dan $N$ subruang komplementer dari suatu ruang Banach~$B$. Jelas bahwa setiap
elemen di $B$ dapat ditulis secara tunggal dalam bentuk $m + n$ untuk suatu $m \in M$ dan $n \in
N$. Definisikan pemetaan $E \colon B \sto B$ dengan $E(m + n) = m$ untuk $m \in M$ dan $n \in N$.
Maka $E$ adalah proyeksi pada $M$ sepanjang~$N$.
\end{prop}

\begin{proof}[\emph{Petunjuk pembuktian}] Gunakan~\ref{prop_cont_idem_op}.  \ns
\end{proof}

\begin{prop}\label{C069817} Jika $E$ adalah proyeksi pada suatu ruang Banach $B$, maka kernel dan jangkauannya merupakan
subruang-subruang komplementer dari~$B$.
\end{prop}

Berikut ini merupakan akibat langsung dari dua proposisi sebelumnya.

\begin{cor}\label{cor_compl_ran_proj} Suatu subruang dari ruang Banach terkomplemen jika dan hanya jika subruang itu
merupakan jangkauan suatu proyeksi.
\end{cor}

\begin{prop} Misalkan $B$ ruang Banach. Jika $E \colon B \sto B$ adalah proyeksi ke suatu subruang
$M$ sepanjang subruang $N$, maka $I - E$ adalah proyeksi ke $N$ sepanjang~$M$.
\end{prop}

\begin{prop}\label{C069824} Jika $M$ dan $N$ merupakan subruang komplementer dari suatu ruang Banach $B$, maka $B/M$ dan
$N$ isomorfik.
\end{prop}

\begin{prop}\label{prop_linv_Bsp} Misalkan $T \colon B \sto C$ pemetaan linear terbatas di antara ruang-ruang Banach.
 \index{kiri!invers!operator ruang Banach}%
Jika $T$ injektif dan $\ran T$ terkomplemen, maka $T$ mempunyai invers kiri.
\end{prop}

\begin{proof}[\emph{Petunjuk pembuktian}] Untuk petunjuk ini kita memperkenalkan notasi (yang agak rumit).
Jika $f\colon X \sto Y$ adalah fungsi di antara dua himpunan, kita definisikan
$f_\bullet \colon X \sto \ran f\colon x \mapsto f(x)$. Artinya, $f_\bullet$ tepat merupakan fungsi yang sama
dengan $f$, kecuali bahwa kodomainnya adalah jangkauan~$f$.

Karena $\ran T$ terkomplemen, terdapat proyeksi $E$ dari $C$ ke~$\ran T$
(menurut~\ref{cor_compl_ran_proj}). Tunjukkan bahwa $T_\bullet$ mempunyai invers kontinu.
Misalkan $S = {T_\bullet}^{-1}E_\bullet$.    \ns
\end{proof}

\begin{prop}\label{prop_rinv_Bsp} Misalkan $S \colon C \sto B$ pemetaan linear terbatas di antara ruang-ruang Banach.
 \index{kanan!invers!operator ruang Banach}%
Jika $S$ surjektif dan $\ker S$ terkomplemen, maka $S$ mempunyai invers kanan.
\end{prop}

\begin{proof}[\emph{Petunjuk pembuktian}] Karena $\ker S$ terkomplemen, terdapat $N \preccurlyeq C$ sedemikian sehingga
$C = N \oplus \ker S$. Tunjukkan bahwa pembatasan $S$ ke $N$ mempunyai invers kontinu dan komposisikan
invers ini dengan pemetaan inklusi dari $N$ ke~$C$.   \ns
\end{proof}

Hasil berikut merupakan konvers dari kedua proposisi~\ref{prop_linv_Bsp} dan~\ref{prop_rinv_Bsp}.

\begin{prop}\label{C069831} Misalkan $T \colon B \sto C$ dan $S\colon C \sto B$ pemetaan linear terbatas di antara ruang-ruang
Banach. Jika $ST = \id B$, maka
 \begin{enumerate}[label=\rm{(\alph*)}]
  \item[(a)] $S$ surjektif;
  \item[(b)] $T$ injektif;
  \item[(c)] $TS$ merupakan proyeksi pada $\ran T$ sepanjang~$\ker S$; dan
  \item[(d)] $C = \ran T \oplus \ker S$.
 \end{enumerate}
\end{prop}

\begin{prop} Jika dalam kategori $\cat{BAN_\infty}$ barisan
 \[ \begin{CD} 0 @>>> A @>j>> B @>f>> C @>>> 0
    \end{CD}\]
 dan
 \[ \begin{CD} 0 @>>> C @>g>> B @>p>> D @>>> 0
    \end{CD}\]
eksak dan jika $fg = \id{C}$, maka $pj$ adalah isomorfisme dari $A$ pada~$D$.
\end{prop}

\begin{proof}[\emph{Petunjuk pembuktian}] Gunakan \ref{C069831} untuk melihat bahwa $B = \ker p \oplus \ran j$. Tinjau pemetaan
$\bigl.p\bigr|_{\ran j}\,j_\bullet$ (notasi untuk $j_\bullet$ seperti dalam petunjuk untuk~\ref{prop_linv_Bsp}). \ns
\end{proof}

\begin{prop}\label{prop_merge_exactCD} Jika dalam kategori $\cat{BAN_\infty}$ diagram
 \begin{equation} \begin{CD} 0 @>>> A @>j>> B @>f>> C @>>> 0  \\
      @.   @VSVV   @VVUV   @.    @.    \\
 0 @>>> A' @>>j'> B' @>>f'> C' @>>> 0
    \end{CD}\end{equation}
dan
 \[ \begin{CD} 0 @>>> C @>g>> B @>p>> D @>>> 0   \\
      @.     @.   @VUVV   @VVTV     @.  \\
  0 @>>> C' @>>g'> B' @>>p'> D' @>>> 0
    \end{CD}\]
komutatif, jika baris-barisnya eksak, jika $fg = \id{C}$ dan jika $f'g' = \id{C'}$, maka diagram
 \[ \begin{CD}  A  @>pj>>  D   \\
       @VSVV   @VVTV  \\
     A'  @>>p'j'> D'
    \end{CD}\]
komutatif dan pemetaan $pj$ serta $p'j'$ merupakan isomorfisme.
\end{prop}

Dalam~\ref{exam_dual_functor} dan~\ref{exam_dualfunctor_exact} kita melihat bahwa dalam kategori ruang Banach
dan pemetaan linear kontinu, funktor dualitas $B \mapsto B^*$, $T \mapsto T^*$ merupakan funktor kontravarian eksak.
Dan dalam~\ref{exam_bar_functor} kita melihat bahwa dalam kategori yang sama pasangan pemetaan $B \mapsto \clo B$,
$T \mapsto \clo T$ merupakan funktor kovarian eksak. Dengan mengomposisikan kedua funktor ini diperoleh dua funktor
kontravarian eksak baru, yang akan kita lambangkan dengan $F$ dan $G$:
  \[ \xymatrix{B\ar[r]^F &{\clo B}^{\,*}}\,,\,\xymatrix{T\ar[r]^F &{\clo T}^{\,*}} \qquad \text{dan} \qquad
           \xymatrix{B\ar[r]^G &\clo{B^*}}\,,\,\xymatrix{T\ar[r]^G &\clo{T^*}}  \]
Wajar untuk menanyakan hubungan antara ${\clo B}^{\,*}$ dan $\clo{B^*}$. Proposisi berikut
memberi tahu kita bukan hanya bahwa kedua ruang ini selalu isomorfik, melainkan bahwa keduanya isomorfik secara alami. Artinya,
 \index{alami!ekuivalensi!antara ${\clo B}^{\,*}$ dan $\clo{B^*}$}%
funktor $F$ dan $G$ ekuivalen secara alami.

\begin{prop} Jika $T \in \ofml B(B,C)$, dengan $B$ dan $C$ ruang Banach, maka terdapat isomorfisme
$\sbsb{\phi}B \colon {\clo B}^{\,*} \sto \clo{B^*}$ dan $\sbsb{\phi}C \colon {\clo C}^{\,*} \sto \clo{C^*}$ sedemikian sehingga
   \[ {\clo T}^{\,*} = {\sbsb{\phi}B}^{\!\!\!-1}\,\, \clo{T^*}\,\,\sbsb{\phi}C\,. \]
\end{prop}

\begin{proof}[\emph{Petunjuk pembuktian}] Karena funktor dualitas eksak (lihat~\ref{exam_dualfunctor_exact}), kita dapat
mengambil dual dari diagram komutatif~\eqref{eqn_exactCD_Bbar} untuk memperoleh diagram komutatif kedua dengan baris-baris
eksak. Sekarang tuliskan diagram komutatif ketiga dengan mengganti setiap kemunculan $B$ dalam diagram~\eqref{eqn_exactCD_Bbar}
dengan~$B^*$. Periksa bahwa ${\sbsb{\tau}B}^{\!\!*}\,\sbsb{\tau}{B^*} = \id{B^*}$ sehingga Anda dapat menggunakan proposisi~\ref{prop_merge_exactCD}
pada dua diagram terakhir dan dengan demikian memperoleh diagram komutatif
     \[ \xy
          \Square/{<-}`{->}`{->}`{<-}/[{\clo B}^{\,*}`{\clo C}^{\,*}`\clo{B^*}`\clo{C^*};{\clo T}^{\,*}`\sbsb{\phi}B`\sbsb{\phi}C`\clo{T^*}]
      \endxy \]
dengan $\sbsb{\phi}B := \sbsb{\pi}{B^*}\,{\sbsb{\pi}B}^{\!\!*}$ suatu isomorfisme.   \ns
\end{proof}

\begin{cor} Suatu ruang Banach refleksif jika dan hanya jika dualnya refleksif.
\end{cor}

\begin{prop}\label{prop_fd_subsp_complm} Setiap subruang berdimensi hingga dari suatu ruang Banach terkomplemen.
\end{prop}

\begin{proof}[\emph{Petunjuk pembuktian}] Misalkan $F$ subruang berdimensi hingga dari suatu ruang Banach $B$ dan misalkan
$\{\vc e^1, \dots, \vc e^n\}$ basis Hamel untuk~$F$. Maka untuk setiap $x \in F$ terdapat skalar-skalar tunggal
$\alpha_1(x)$, \dots , $\alpha_n(x)$ sedemikian sehingga $x = \sum_{k=1}^n \alpha_k(x)\vc e^k$. Untuk setiap $1 \le k \le n$
verifikasi bahwa $\alpha_k$ merupakan fungsional linear terbatas, yang dapat diperluas ke seluruh~$B$. Misalkan $N$ adalah
irisan kernel-kernel dari perluasan ini dan tunjukkan bahwa $B = F \oplus N$.  \ns
\end{proof}

\begin{defn} Ingat bahwa dalam~\ref{000082} kita mendefinisikan \emph{kodimensi} suatu subruang dari ruang vektor
sebagai dimensi komplemen ruang vektornya (yang selalu ada menurut proposisi~\ref{0000824}). Kita menggunakan
istilah yang sama untuk ruang Banach:
 \index{kodimensi}%
\df{kodimensi} dari sebarang \emph{subruang vektor} suatu ruang Banach adalah dimensi \emph{komplemen ruang
vektornya}.
\end{defn}

Untuk korolar berikut (dari proposisi~\ref{prop_fd_subsp_complm}), ingatlah
konvensi~\ref{conv_subsp_Bsp} tentang penggunaan kata ``subruang'' dalam konteks ruang Banach.

\begin{cor} Dalam suatu ruang Banach, setiap subruang berkodimensi hingga terkomplemen.
\end{cor}





















Misalkan $M$ subruang tak nol dari suatu ruang Banach~$B$. Tinjau dua sifat berikut dari pasangan $B$ dan~~$M$.

\textbf{Sifat P:} \emph{Terdapat proyeksi bernorma $1$ dari $B$ ke $M$.}

\textbf{Sifat E:} \emph{Untuk setiap ruang Banach $C$, setiap pemetaan linear terbatas dari $M$ ke $C$ dapat diperluas menjadi pemetaan
linear terbatas dari $B$ ke $C$ tanpa memperbesar normanya.}

\begin{prop} Untuk subruang tak nol $M$ dari suatu ruang Banach $B$, sifat \textbf{P} dan \textbf{E} ekuivalen.
\end{prop}

\begin{proof}[\emph{Petunjuk pembuktian}] Untuk menunjukkan bahwa \textbf{E} mengakibatkan \textbf{P}, mulailah dengan memperluas pemetaan identitas
pada $M$ ke seluruh~$B$.   \ns
\end{proof}

















\section{Prinsip Keterbatasan Seragam}
\begin{defn} Misalkan $S$ suatu himpunan dan $V$ ruang linear bernorma. Suatu keluarga $\fml F$ dari
fungsi-fungsi dari $S$ ke $V$ disebut
 \index{titik demi titik!terbatas}%
 \index{terbatas!titik demi titik}%
\df{terbatas titik demi titik} jika untuk setiap $x \in S$ terdapat konstanta $M_x > 0$ sedemikian sehingga
$\norm{f(x)} \le M_x$ untuk setiap $f \in \fml F$.
\end{defn}

\begin{defn} Misalkan $V$ dan $W$ ruang linear bernorma. Suatu keluarga $\fml T$ dari pemetaan linear terbatas
dari $V$ ke $W$ disebut
 \index{seragam!terbatas}%
 \index{terbatas!seragam}%
\df{terbatas seragam} jika terdapat $M > 0$ sedemikian sehingga $\norm T \le M$ untuk setiap $T \in
\fml T$.
\end{defn}

\begin{exer}\label{C070114} Misalkan $V$ dan $W$ ruang linear bernorma. Jika suatu keluarga $\fml T$ dari
pemetaan linear terbatas dari $V$ ke $W$ terbatas seragam, maka keluarga itu terbatas titik demi titik.
\end{exer}

\emph{Prinsip keterbatasan seragam} adalah pernyataan bahwa konvers
dari~\ref{C070114} berlaku apabila $V$ lengkap. Untuk membuktikannya kita akan menggunakan
analog ``lokal'' dari pernyataan ini untuk fungsi-fungsi kontinu pada ruang metrik lengkap. Secara kasar, analog ini mengatakan bahwa jika
suatu keluarga fungsi bernilai real pada ruang metrik lengkap terbatas titik demi titik, maka keluarga itu terbatas
seragam pada suatu subhimpunan terbuka tak kosong dari ruang tersebut.

\begin{prop}\label{C070117} Misalkan $M$ ruang metrik lengkap yang tak kosong dan $\fml F \subseteq \fml C(M,\R)$.
Jika untuk setiap $x \in M$ terdapat konstanta $N_x > 0$ sedemikian sehingga $\abs{f(x)} \le N_x$ untuk
setiap $f \in \fml F$, maka terdapat himpunan terbuka tak kosong $U$ di $M$ dan bilangan $N > 0$ sedemikian
sehingga $\abs{f(u)} \le N$ untuk setiap $f \in \fml F$ dan setiap $u \in U$.
\end{prop}

\begin{proof}[\emph{Petunjuk pembuktian}] Untuk setiap bilangan asli $k$, misalkan $A_k =
\bigcap\{f^{\gets}\bigl(\,[-k,k]\,\bigr)\colon f \in \fml F\}$. Gunakan versi yang sama dari
\emph{teorema kategori Baire} yang kita gunakan dalam membuktikan \emph{teorema pemetaan terbuka} (\ref{C069414})
untuk menyimpulkan bahwa untuk sedikitnya satu $n \in \N$, himpunan $A_n$ mempunyai interior
tak kosong. Misalkan $U = \intr{A_n}$.   \ns
\end{proof}

\begin{thm}[Prinsip keterbatasan seragam]\label{thm_PUB} Misalkan $B$ ruang Banach dan $W$ ruang linear bernorma.
 \index{prinsip!keterbatasan seragam}%
 \index{keterbatasan!prinsip seragam}%
 \index{seragam!keterbatasan!prinsip}%
Jika suatu keluarga $\ofml T$ dari pemetaan linear terbatas dari $B$ ke $W$ terbatas titik demi
titik, maka keluarga itu terbatas seragam.
\end{thm}

\begin{proof}[\emph{Petunjuk pembuktian}] Untuk setiap $T \in \ofml T$, definisikan
  \[ \sbsb fT\colon B \sto \R\colon x \mapsto \norm{Tx}. \]
Gunakan proposisi sebelumnya untuk menunjukkan bahwa terdapat subhimpunan terbuka tak kosong $U$ dari $B$
dan konstanta $M > 0$ sedemikian sehingga $\sbsb fT(x) \le M$ untuk setiap $T \in \ofml T$ dan setiap
$x \in U$. Pilih titik $a \in B$ dan bilangan $r > 0$ sedemikian sehingga $\clo{B_r(a)} \subseteq
U$. Kemudian verifikasi bahwa
  \[ \norm{Ty} \le r^{-1}\bigl(\, \sbsb fT(a + ry) + \sbsb fT(a)\,\bigr) \le 2Mr^{-1} \]
untuk setiap $T \in \ofml T$ dan setiap $y \in B$ sedemikian sehingga $\norm y \le 1$. \ns
\end{proof}

\begin{thm}[Banach-Steinhaus] Misalkan $B$ ruang Banach, $W$ ruang linear bernorma,
 \index{teorema Banach-Steinhaus}%
dan $(T_n)$ barisan pemetaan linear terbatas dari $B$ ke~$W$. Jika limit titik demi titik
$\lim_{n \sto \infty} T_nx$ ada untuk setiap $x$ di~$B$, maka pemetaan $S\colon B \sto
W\colon x \mapsto \lim_{n \sto \infty} T_nx$ adalah pemetaan linear terbatas.
\end{thm}

\begin{proof}[\emph{Petunjuk pembuktian}] Untuk setiap $x \in B$, barisan $\bigl(\norm{T_nx}\bigr)_{n=1}^\infty$
konvergen, dan karena itu terbatas. Gunakan~\ref{thm_PUB}.   \ns
\end{proof}

\begin{defn} Suatu subhimpunan $A$ dari ruang linear bernorma $V$ disebut
 \index{terbatas!secara lemah!himpunan dalam ruang linear bernorma}%
 \index{lemah!keterbatasan!dalam ruang linear bernorma}%
 \index{w@$w$-terbatas}%
\df{terbatas secara lemah} ($w$-terbatas) jika $f^{\sto}(A)$ adalah subhimpunan terbatas dari medan skalar
$V$ untuk setiap $f \in V^*$.
\end{defn}

\begin{prop} Suatu subhimpunan $A$ dari ruang Hilbert $H$ terbatas secara lemah jika dan hanya jika untuk setiap $y \in H$
 \index{lemah!keterbatasan!dalam ruang Hilbert}%
himpunan $\{\langle a,y \rangle \colon a \in A\}$ terbatas.
\end{prop}

Proposisi sebelumnya kurang lebih sudah jelas. Proposisi berikutnya merupakan konsekuensi penting dari
\emph{prinsip keterbatasan seragam}.

\begin{prop}\label{C070131} Suatu subhimpunan dari ruang linear bernorma terbatas secara lemah jika dan hanya jika
subhimpunan itu terbatas.
\end{prop}

\begin{proof}[\emph{Petunjuk pembuktian}] Jika $A$ subhimpunan yang terbatas secara lemah dari ruang linear bernorma $V$
dan $f \in V^*$, apa yang dapat Anda katakan tentang $\sup\{\abs{a^{**}(f)}\colon a \in A\}$?  \ns
\end{proof}

\begin{exer} Sahabat baik Anda, Fred R. Dimm, kembali memerlukan bantuan. Kali ini ia mencemaskan
pernyataan (lihat~\ref{C070131}) bahwa suatu subhimpunan ruang linear bernorma terbatas jika
dan hanya jika subhimpunan itu terbatas secara lemah. Ia meninjau barisan $(a^n)$ dalam ruang Hilbert
$l^2$ yang didefinisikan oleh $a^n = \sqrt{n}\,\vc e^n$ untuk setiap $n \in \N$, dengan $\{\vc
e^n\colon n \in \N\}$ basis ortonormal biasa untuk~$l_2$ (lihat
contoh~\ref{C063149}). Ia melihat bahwa barisan itu jelas tidak terbatas (dengan topologi biasa
pada~$l^2$). Namun, baginya barisan itu tampak terbatas secara lemah. Jika tidak, menurut argumennya, maka
berdasarkan \emph{teorema Riesz-Fr\'echet} \ref{000342} akan terdapat barisan $x$ dalam
$l^2$ sedemikian sehingga himpunan $\{\abs{\langle a^n,x \rangle}\colon n \in\N\}$ tidak terbatas. Baginya ini tampak
tidak mungkin terjadi karena barisan semacam itu harus menurun lebih lambat daripada
barisan $\bigl(n^{-1/2}\bigr)$, yang sudah menurun terlalu lambat untuk menjadi anggota~$l^2$.
Tenangkan Fred dengan menemukan barisan yang sesuai. \emph{Petunjuk:} Gunakan barisan yang
memiliki banyak nol.
\end{exer}

\begin{defn} Suatu barisan $(x_n)$ dalam ruang linear bernorma $V$ disebut
 \index{lemah!barisan Cauchy}%
 \index{Cauchy!barisan!lemah}%
 \index{barisan!Cauchy lemah}%
\df{Cauchy secara lemah} jika barisan $(f(x_n))$ bersifat Cauchy untuk setiap $f$ dalam~$V^*$. Barisan
tersebut
 \index{lemah!konvergensi!barisan dalam ruang linear bernorma}%
 \index{konvergensi!lemah}%
 \index{barisan!konvergensi lemah dari suatu}%
\df{konvergen secara lemah} ke titik $a \in V$ jika $f(x_n) \sto f(a)$ untuk setiap $f \in V^*$
(yakni, jika barisan itu konvergen dalam topologi lemah yang diinduksi oleh anggota-anggota~$V^*$). Dalam hal ini
kita tulis
 \index{<arrowtow@$x_n \to^w a$ (konvergensi sekuensial lemah)}%
$x_n \to^w a$. Ruang $V$ disebut
 \index{lengkap!secara lemah}%
 \index{lemah!kelengkapan}%
\df{lengkap sekuensial secara lemah} jika setiap barisan Cauchy secara lemah di $V$ konvergen secara lemah.
\end{defn}

\begin{prop} Dalam ruang linear bernorma, setiap barisan Cauchy secara lemah terbatas.
\end{prop}

\begin{prop} Setiap ruang Banach refleksif lengkap sekuensial secara lemah.
\end{prop}

\begin{exam} Misalkan $f_n(t) = \left\{
                             \begin{array}{ll}
                               1 - nt, & \hbox{jika $0 \le t \le \frac1n$;} \\
                               0, & \hbox{jika $\frac1n < t \le 1$.}
                             \end{array}
                           \right.$ untuk setiap $n \in \N$. Barisan fungsi $(f_n)$ menunjukkan bahwa ruang Banach
$\fml C([0,1])$ tidak lengkap sekuensial secara lemah.
\end{exam}

\begin{cor} Ruang Banach $\fml C([0,1])$ tidak refleksif.
\end{cor}

\begin{prop} Misalkan $\fml E$ basis ortonormal untuk suatu ruang Hilbert~$H$ dan $(x_n)$ suatu
barisan di~$H$. Maka $x_n \to^w \vc 0$ jika dan hanya jika $(x_n)$ terbatas dan $\langle
x_n,e \rangle \sto 0$ untuk setiap $e \in \fml E$.
\end{prop}

\begin{defn} Misalkan $H$ ruang Hilbert. Suatu jaring $(T_\lambda)$ dalam $\ofml B(H)$
 \index{lemah!topologi operator!konvergensi dalam}%
 \index{operator!topologi!lemah}%
 \index{topologi!operator lemah}%
 \index{lemah!konvergensi!operator ruang Hilbert}%
\df{konvergen secara lemah} (atau \df{konvergen dalam topologi operator lemah}) ke $S \in \ofml B(H)$ jika
  \[ \langle T_\lambda x,y \rangle \sto \langle Sx,y \rangle \]
untuk setiap $x$, $y \in H$. Dalam hal ini kita tulis
$T_\lambda~\to^{\textbf{WOT}}~S$.
 \index{<arrowtowot@$T_\lambda~\to^{\textbf{WOT}}~S$}%

Suatu keluarga $\ofml T \subseteq \ofml B(H)$ disebut
 \index{lemah!keterbatasan!keluarga operator ruang Hilbert}%
 \index{terbatas!secara lemah}
 \index{terbatas!dalam topologi operator lemah}%
 \index{lemah!topologi operator!keterbatasan dalam}%
 \index{WOT@\textbf{WOT} (topologi operator lemah)}%
\df{terbatas secara lemah} (atau \df{terbatas dalam topologi operator lemah}) jika untuk setiap $x$, $y \in H$
terdapat konstanta positif $\alpha_{x,y}$ sedemikian sehingga
  \[\abs{\langle Tx,y \rangle} \le \alpha_{x,y}\]
untuk setiap $T \in \ofml T$.
\end{defn}

Sayangnya definisi sebelumnya dapat menjadi sumber kebingungan. Karena $\ofml B(H)$ adalah
ruang linear bernorma, ruang itu sudah mempunyai topologi ``lemah'' (lihat~\ref{C067434}). Lalu, bagaimana menangani dua
topologi berbeda pada $\ofml B(H)$ yang mempunyai nama sama? Sebagian penulis menyebut topologi yang baru saja didefinisikan sebagai
\emph{topologi operator lemah}, sehingga mencegah timbulnya kebingungan. Penulis lain berpendapat bahwa
topologi lemah dalam pengertian ruang linear bernorma begitu jarang berguna dalam teori operator pada ruang Hilbert
sehingga, kecuali diberi penjelas khusus, frasa \emph{topologi lemah} dalam konteks operator ruang Hilbert
seharusnya dipahami sebagai topologi yang baru saja didefinisikan. (Lihat, misalnya, komentar Halmos dalam dua
paragraf sebelum Soal 108 di~\cite{Halmos:1982}.) Demikian pula, kita sering melihat frasa \emph{suatu
himpunan terbatas yang terdiri atas operator-operator ruang Hilbert.} Kecuali diberi penjelas lain, yang dimaksud adalah bahwa himpunan
tersebut \emph{terbatas seragam}. Bagaimanapun juga, berhati-hatilah terhadap frasa-frasa ini ketika membaca.

Berikut ini konsekuensi penting dari \emph{prinsip keterbatasan seragam}~\ref{thm_PUB}.

\begin{prop} Setiap keluarga operator ruang Hilbert yang terbatas dalam topologi operator
lemah adalah terbatas seragam.
\end{prop}

\begin{proof}[\emph{Petunjuk pembuktian}] Misalkan $H$ ruang Hilbert dan $\ofml T$ keluarga operator pada~$H$
yang terbatas terhadap topologi operator lemah. Gunakan proposisi~\ref{C070131} untuk menunjukkan, bagi
$y \in H$ yang tetap, bahwa $\{T^*y\colon T \in \ofml T\}$ adalah subhimpunan terbatas dari~$H$. Gunakan hal ini (dan~\ref{C070131}
sekali lagi) untuk menunjukkan bahwa $\{Tx\colon \text{$T \in \ofml T$, $x \in H$, dan $\norm x \le 1$}\}$ terbatas di~$H$. \ns
\end{proof}

Seperti dinyatakan di atas, banyak penulis akan menyingkat pernyataan proposisi sebelumnya menjadi:
\emph{Setiap keluarga operator ruang Hilbert yang terbatas secara lemah adalah terbatas.}

\begin{defn}
Suatu jaring $(T_\lambda)$ dalam $\ofml B(H)$
 \index{kuat!topologi operator!konvergensi dalam}%
 \index{operator!topologi!kuat}%
 \index{topologi!operator kuat}%
 \index{kuat!konvergensi!operator ruang Hilbert}%
 \index{SOT@\textbf{SOT} (topologi operator kuat)}%
\df{konvergen secara kuat} (atau \df{konvergen dalam topologi operator kuat}) ke $S \in \ofml B(H)$ jika
  \[ T_\lambda x \sto Sx \]
untuk setiap $x \in H$. Dalam hal ini kita tulis $T_\lambda~\to^{\textbf{SOT}}~S$.
 \index{<arrowtosot@$T_\lambda~\to^{\textbf{SOT}}~S$}%

Konvergensi norma biasa dari operator-operator dalam $\ofml B(H)$ sering disebut
 \index{seragam!topologi operator!konvergensi dalam}%
 \index{operator!topologi!seragam}%
 \index{topologi!operator seragam}%
 \index{seragam!konvergensi!operator ruang Hilbert}%
\df{konvergensi seragam} (atau \df{konvergensi dalam topologi operator seragam}). Artinya, kita tulis
 \index{<arrowtouniform@$T_\lambda~\sto~S$}%
$T_\lambda \sto S$ jika $\norm{S - T_\lambda}\sto 0$.
\end{defn}

\begin{prop} Misalkan $H$ ruang Hilbert, $(T_n)$ barisan dalam $\ofml B(H)$,
dan $B \in \ofml B(H)$. Maka implikasi-implikasi berikut berlaku:
   \[ T_n \sto B \quad \implies \quad T_n~\to^{\textbf{SOT}}~B
                                \quad \implies \quad T_n~\to^{\textbf{WOT}}~B. \]
\end{prop}

\begin{prop} Misalkan $(S_n)$ dan $(T_n)$ barisan operator pada suatu ruang Hilbert~$H$. \\
Jika $S_n \to^{\textbf{WOT}} A$ dan $T_n \to^{\textbf{SOT}} B$, maka $S_nT_n \to^{\textbf{WOT}} AB$.
\end{prop}

\begin{proof}[\emph{Petunjuk pembuktian}] Tuliskan $S_nT_n - AB$ sebagai $S_nT_n - S_nB + S_nB - AB$.  \ns
\end{proof}

\begin{exam} Dalam proposisi sebelumnya, hipotesis $T_n \to^{\textbf{SOT}} B$ tidak dapat
diganti dengan $T_n~\to^{\textbf{WOT}}~B$.
\end{exam}

\begin{proof}[\emph{Petunjuk pembuktian}] Tinjau pangkat-pangkat operator geser unilateral.  \ns
\end{proof}






\endinput
